The glymphatic hypothesis predicts that reduced slow-wave sleep permits amyloid-beta to
accumulate, and this prediction now drives trials that target N3 sleep duration as a route to
slowing Alzheimer's pathology. No meta-analysis has ever tested it: prior syntheses pooled
sleep quality, sleep duration or apnea severity, never a slow-wave-specific exposure against
amyloid PET. We pool 40 effect sizes from 17 studies and 12 cohort families, keeping N3
\emph{stage} amount strictly separate from \emph{spectral} slow-wave activity, aggregating
dependent effects to one synthetic effect per cohort family, and recovering correlations from
reported regression coefficients so that large cohorts are not silently dropped. The pooled
N3-stage/amyloid association is $r = +0.026$ (95\% CI $-0.153$ to $+0.203$, $p = 0.68$,
$I^2 = 0\%$), equivalence-tested tightly enough to reject a true correlation as large as
$r = 0.20$ ($p_{\mathrm{TOST}} = 0.026$). It is weaker, not stronger, than total sleep time
($r = +0.163$) and sleep efficiency ($r = +0.090$); all three preregistered specificity
contrasts point opposite to the hypothesis. We introduce a measurement-reliability null model
built from empirical single-night ICCs and show it predicts N3 should look $1.86\times$
\emph{stronger} than total sleep time; the observed ratio is $0.16\times$. Converted to the
hypothesis's own units, a 10-percentage-point reduction in N3 gives
$\Delta\mathrm{SUVR} = +0.005$, $1.8\%$ of the amyloid-negative/positive separation. The
positive spectral literature collapses from $r = +0.112$ to $r = -0.025$ when one cohort family
is removed.
